\subsection{\texorpdfstring{\(\mathfrak{sl}_2\) 三元组在半单李代数中}{\textbackslash mathfrak\{sl\}\_2 三元组在半单李代数中}}

引理 5。令 V 为有限维向量空间，A、B 为 V 的自同态。假设 A 是幂零的且 \([A, [A, B]] = 0\)。则 AB 是幂零的。

置 C = {[}A, B{]}。由于 {[}A, C{]} = 0，

\[
[ \mathrm{A}, \mathrm{BC} ^ {p} ] = [ \mathrm{A}, \mathrm{B} ] \mathrm{C} ^ {p} = \mathrm{C} ^ {p + 1}
\]

对每个整数 \(p \geq 0\) 都成立。因此，\(\mathrm{Tr}(\mathrm{C}^{p}) = 0\) 对 \(p \geq 1\) 成立，这证明 C 是幂零的（《代数》第 VII 章第3节第5号命题13的推论4）。现在令 \(\bar{k}\) 为 k 的代数闭包，令 \(\lambda \in \bar{k}\)、\(x \in V \otimes_{k} \bar{k}\) 满足 \(\lambda x\)、\(x \neq 0\)。关系 \([[B, A], A] = 0\) 表明，对每个整数 \([B, A^{p}] = p[B, A]A^{p-1}\)，都有 \(p \geq 0\)。令 r 为满足 \(A^{r}x = 0\) 的最小整数。于是

\[
\lambda \mathrm{A} ^ {r - 1} x = \mathrm{A} ^ {r - 1} \mathrm{AB} x = \mathrm{A} ^ {r} \mathrm{B} x = \mathrm{BA} ^ {r} x - [ \mathrm{B}, \mathrm{A} ^ {r} ] x = - r [ \mathrm{B}, \mathrm{A} ] \mathrm{A} ^ {r - 1} x.
\]

由于 {[}B, A{]} 是幂零的且 \(A^{r-1}x \neq 0\)，可知 \(\lambda = 0\)。因此 AB 的所有特征值都为零，故引理得证。

引理 6。令 \(h, x \in \mathfrak{g}\) 满足 \([h, x] = 2x\)，且 \(h \in (\operatorname{ad} x)(\mathfrak{g})\)。则存在 \(y \in \mathfrak{g}\)，使得 \((x, h, y)\) 或为 \((0, 0, 0)\)，或为一个 \(\mathfrak{sl}_2\) 三元组。

令 \(\mathfrak{g}'\) 为可解李代数 \(kh + kx\)。由于 \(x \in [\mathfrak{g}', \mathfrak{g}']\)，\(\operatorname{ad}_{\mathfrak{g}} x\) 是幂零的（第 I 章第5节第3号定理1）；令 \(\mathfrak{n}\) 为其核。由于 {[}ad \(h\) , ad \(x] = 2\) ad \(x\)，有 (ad \(h\) ) \(\mathfrak{n} \subset \mathfrak{n}\)。令 \(z \in \mathfrak{g}\) 满足 \(h = -[x, z]\)。对任意整数 \(n \geq 0\)，置 \(\mathrm{M}_n = (\mathrm{ad}~x)^n\mathfrak{g}\)。若 \(n > 0\)，则有（§1，第1号，引理1）

\[
[ \mathrm{ad} z, (\mathrm{ad} x) ^ {n} ] = n ((\mathrm{ad} h) - n + 1) (\mathrm{ad} x) ^ {n - 1}
\]

所以，若 \(u \in M_{n-1}\)，

\[
n ((\operatorname{ad} h) - n + 1) u \in (\operatorname{ad} z) (\operatorname{ad} x) u + \mathrm{M} _ {n}.
\]

由于 (ad \(h)\mathfrak{n}\subset \mathfrak{n}\)，可得

\[
((\operatorname{ad} h) - n + 1) (\mathfrak {n} \cap \mathrm{M} _ {n - 1}) \subset \mathfrak {n} \cap \mathrm{M} _ {n}.
\]

由于 ad x 是幂零的，\(M_{n} = 0\) 对充分大的 n 成立。~因此，ad \(h|n\) 的特征值都是 \(\geq 0\) 的整数。于是 ad \(h + 2\) 在 n 上的限制可逆。

现在 \([h,z] + 2z\in \mathfrak{n}\)，因为

\[
\begin{array}{r} [ x, [ h, z ] + 2 z ] = [ [ x, h ], z ] + [ h, [ x, z ] ] + 2 [ x, z ] \\ = [ - 2 x, z ] + [ h, - h ] + 2 [ x, z ] = 0. \end{array}
\]

因此存在 \(z' \in \mathfrak{n}\)，使 \([h, z'] + 2z' = [h, z] + 2z\)，即置 \([h, y] = -2y\) 后有 \(y = z - z'\)。由于 \([x, y] = [x, z] = -h\)，证明完成。

命题 2（Jacobson–Morozov）。假设 \(\mathfrak{g}\) 是半单的。令 \(x\) 为 \(\mathfrak{g}\) 的非零幂零元素。存在 \(h, y \in \mathfrak{g}\)，使得 \((x, h, y)\) 是一个 \(\mathfrak{sl}_2\) 三元组。

令 \(\mathfrak{n}=\operatorname{Ker}(\operatorname{ad}x)^{2}\)。若 \(z\in\mathfrak{n}\)，则 {[}ad x, {[}ad x, ad z{]}{]} = ad({[}x, {[}x,z{]}{]}) = 0。由引理5，ad \(x\circ ad\) z 是幂零的，所以 \(\operatorname{Tr}(\operatorname{ad}x\circ\operatorname{ad}z)=0\)。这表明 x 关于 g 的 Killing 型 \(\Phi\) 与 n 正交。由于

\[
\Phi ((\operatorname{ad} x) ^ {2} y, y ^ {\prime}) = \Phi (y, (\operatorname{ad} x) ^ {2} y ^ {\prime})
\]

对所有 \(y, y' \in g\) 都成立，并且由于 \(\Phi\) 非退化，n 的正交补是 \((\operatorname{ad} x)^{2}\) 的像。因此存在 \(y' \in g\)，使 \(x = (\operatorname{ad} x)^{2}y'\)。置

\[
h = - 2 [ x, y ^ {\prime} ];
\]

我们有 \([h,x]=2x\) 且 \(h\in(\operatorname{ad}x)(\mathfrak{g})\)。现在只需应用引理6。

推论。假设 \(\mathfrak{g}\) 是半单的。令 \(\mathrm{G}\) 为包含 \(\mathfrak{g}\) 的 \(\operatorname{Aut}_e(\mathfrak{g})\) 的自同构群。将 g 中任意 \(\mathfrak{sl}_2\) 三元组 \((x,h,y)\) 与幂零元素 x 关联的映射在取商后，定义了从 \(sl_{2}\) 三元组的 G-共轭类到非零幂零元素的 G-共轭类的一个双射。

这由命题1和2 得出。

引理 7。令 K 为至少含有 4 个元素的交换域。令 G 为矩阵 \(\begin{pmatrix}\alpha&\beta\\0&\alpha^{-1}\end{pmatrix}\) 所成的群，其中 \(\alpha\in K^{*},\beta\in K\)。令 \(G'\) 为满足 \(\alpha=1\) 的这类矩阵所成的群。则 \(G'=(G,G)\)。

若 \(\alpha, \alpha' \in \mathrm{K}^*\) 且 \(\beta, \beta' \in \mathrm{K}\)，

\[
\begin{array}{c} \left( \begin{array}{c c} \alpha & \beta \\ 0 & \alpha^ {- 1} \end{array} \right) \left( \begin{array}{c c} \alpha^ {\prime} & \beta^ {\prime} \\ 0 & \alpha^ {- 1} \end{array} \right) \left( \begin{array}{c c} \alpha & \beta \\ 0 & \alpha^ {- 1} \end{array} \right) ^ {- 1} \left( \begin{array}{c c} \alpha^ {\prime} & \beta^ {\prime} \\ 0 & \alpha^ {- 1} \end{array} \right) ^ {- 1} \\ = \left( \begin{array}{c c} 1 & - \alpha^ {\prime} \beta^ {\prime} - \alpha \beta \alpha^ {2} + \alpha^ {2} \alpha^ {\prime} \beta^ {\prime} + \alpha \beta \\ 0 & 1 \end{array} \right). \end{array}
\]

特别地，

\[
\left( \begin{array}{c c} 1 & \beta \\ 0 & 1 \end{array} \right) \left( \begin{array}{c c} \alpha^ {\prime} & 0 \\ 0 & \alpha^ {- 1} \end{array} \right) \left( \begin{array}{c c} 1 & \beta \\ 0 & 1 \end{array} \right) ^ {- 1} \left( \begin{array}{c c} \alpha^ {\prime} & 0 \\ 0 & \alpha^ {- 1} \end{array} \right) ^ {- 1} = \left( \begin{array}{c c} 1 & \beta (1 - \alpha^ {2}) \\ 0 & 1 \end{array} \right).
\]

但存在 \(\alpha_0' \in \mathrm{K}^*\)，使 \(\alpha_0' \neq 1\) 且 \(\alpha_0' \neq -1\)，于是 \(k.(1 - \alpha_0'^2) = k\)，故引理得证。

命题 3。假设 \(\mathfrak{g}\) 是半单的。群 \(\mathrm{Aut}_e(\mathfrak{g})\) 等于其换位子群。若 \(\mathfrak{g}\) 可分裂，则 \(\mathrm{Aut}_e(\mathfrak{g})\) 是 \(\mathrm{Aut}_0(\mathfrak{g})\) 的换位子群。

令 \(x\) 为 \(\mathfrak{g}\) 的非零幂零元素。取 \(h, y \in \mathfrak{g}\)\((x, h, y)\)，使  为一个 \(\mathfrak{sl}_2\) 三元组（命题2）。由  生成的  的子代数 \(\mathfrak{s}\) 可识别为 。令  为  在 \(\mathfrak{g}\) 上的表示 \((x, h, y)\)\(\mathfrak{sl}(2, k)\)，并令  为与 \(\rho\) 相容的 \(z \mapsto \operatorname{ad}_{\mathfrak{g}} z\)\(\mathfrak{s} = \mathfrak{sl}(2, k)\)\(\mathfrak{g}\) 的表示（§1，第4号）。\(\pi\) 的像由 \(\mathbf{SL}(2, k)\)\(\rho\)\(\S 1\)\(\pi\) 的 \(\exp(t\operatorname{ad}_{\mathfrak{g}} x)\) 和 \(\exp(t\operatorname{ad}_{\mathfrak{g}} y)\) 生成（《代数》第 III 章第8节第9号命题17），因而包含于 \(t \in k\)\(\S 8\)\(\operatorname{Aut}_e(\mathfrak{g})\) 中。由于 \(\mathbf{SL}(2, k)\) 等于其换位子群（引理7及同上），\(\exp(\operatorname{ad}_{\mathfrak{g}} x)\) 属于 \(\operatorname{Aut}_e(\mathfrak{g})\) 的换位子群 G。因此 \(\operatorname{Aut}_e(\mathfrak{g})\) 等于 G。现在假设 \(\mathfrak{g}\) 可分裂。由于 \(\operatorname{Aut}_0(\mathfrak{g}) / \operatorname{Aut}_e(\mathfrak{g})\) 是交换的（§5，第3号，注3），上面的论证说明 \(\S 5\)\(\operatorname{Aut}_0(\mathfrak{g})\) 的换位子群是 \(\operatorname{Aut}_e(\mathfrak{g})\)。

